\documentclass[11pt,a4paper]{article}

\usepackage{fontspec}
\usepackage{polyglossia}
\setdefaultlanguage{russian}
\setotherlanguage{english}
\setmainfont[Ligatures=TeX]{STIX Two Text}
\setmonofont{PT Mono}[Scale=0.92]
\newfontfamily\cyrillicfont[Ligatures=TeX]{STIX Two Text}
\newfontfamily\cyrillicfonttt{PT Mono}[Scale=0.92]
% STIX Two Text doesn't have all symbols — fall back to system font
% for arrows and box-drawing chars used in architecture diagrams.
\usepackage{newunicodechar}
\newfontfamily\fallbackfont{Apple Symbols}
\newunicodechar{→}{{\fallbackfont →}}
\usepackage[a4paper,margin=2.4cm]{geometry}
\usepackage{tabularx}
\usepackage{booktabs}
\usepackage{titling}
\usepackage{titlesec}
\usepackage[hidelinks]{hyperref}
\usepackage{xcolor}
\usepackage{enumitem}
\usepackage{listings}
\usepackage{tcolorbox}
\usepackage{graphicx}
\usepackage{fancyhdr}

\definecolor{accent}{HTML}{7c3aed}
\definecolor{ok}{HTML}{10b981}
\definecolor{warn}{HTML}{f59e0b}
\definecolor{err}{HTML}{ef4444}
\definecolor{codecomment}{HTML}{8b8b95}
\definecolor{codecmd}{HTML}{c2410c}
\definecolor{codestr}{HTML}{0369a1}
\definecolor{boxbg}{HTML}{f5f3ff}
\definecolor{boxborder}{HTML}{c4b5fd}

\hypersetup{
  colorlinks=false,
  pdftitle={EclipVPN — Руководство пользователя},
  pdfauthor={EclipVPN},
  pdfsubject={Установка и настройка EclipVPN на OpenWrt},
}

\titleformat{\section}{\Large\bfseries\color{accent}}{\thesection}{1em}{}
\titleformat{\subsection}{\large\bfseries}{\thesubsection}{1em}{}
\titlespacing*{\section}{0pt}{1.5em}{0.6em}
\titlespacing*{\subsection}{0pt}{1.0em}{0.4em}

\setlist[itemize]{leftmargin=*,topsep=2pt,itemsep=2pt}
\setlist[enumerate]{leftmargin=*,topsep=2pt,itemsep=2pt}

\lstdefinelanguage{shell}{
  morekeywords={eclipvpn,uci,opkg,nft,ip,jq,ssh,curl,cat,chmod,sleep,grep,tail,head,wc,for,do,done,killall},
  morecomment=[l]\#,
  morestring=[b]",
  morestring=[b]',
  basicstyle=\ttfamily\small,
  keywordstyle=\color{codecmd}\bfseries,
  commentstyle=\color{codecomment}\itshape,
  stringstyle=\color{codestr},
  breaklines=true,
  showstringspaces=false,
  keepspaces=true,
  columns=fullflexible,
  frame=none,
  xleftmargin=0pt,
  aboveskip=4pt,
  belowskip=4pt,
  extendedchars=true,
  inputencoding=utf8,
}
\lstset{language=shell}

\newtcolorbox{notebox}{
  colback=boxbg,
  colframe=boxborder,
  boxrule=0.6pt,
  arc=3pt,
  left=10pt,right=10pt,top=6pt,bottom=6pt,
  fontupper=\small
}

\pagestyle{fancy}
\fancyhf{}
\fancyfoot[C]{\thepage}
\fancyhead[L]{\small\color{gray}EclipVPN}
\fancyhead[R]{\small\color{gray}Руководство пользователя}
\renewcommand{\headrulewidth}{0pt}

\title{\vspace{2cm}\Huge\textbf{EclipVPN}\\[0.4em]\large\mdseries Руководство по установке и настройке}
\author{}
\date{\today}

\begin{document}
\maketitle
\thispagestyle{empty}
\vspace{1cm}

\begin{center}
\small Лёгкий VPN-менеджер для OpenWrt \\
VLESS+XHTTP+Reality · sing-box tproxy · Веб-интерфейс
\end{center}

\vspace{2cm}
\tableofcontents
\thispagestyle{empty}
\newpage

\section{Введение}

EclipVPN направляет трафик домашней сети через VPN-сервер с поддержкой VLESS+XHTTP+Reality (xray-core). Российские сайты идут напрямую благодаря geosite-ru/geoip-ru. Управление через веб-интерфейс на самом роутере.

\subsection{Основные возможности}

\begin{itemize}
\item \textbf{Два режима маршрутизации}: всё через VPN (с прямым доступом к RU) или только заблокированные сайты
\item \textbf{VLESS+XHTTP+Reality} через xray-core 26.4.17+
\item \textbf{Защита DNS}: перехват DNS, блокировка DoT/QUIC от подмены провайдером
\item \textbf{Без утечек}: при перезапуске nftables и dnsmasq остаются на месте
\item \textbf{Авто-переключение}: watchdog заменяет упавший сервер (без тарифицируемых операторов)
\item \textbf{Маршрутизация по устройствам}: VPN или Direct для каждого клиента
\item \textbf{Скорость в реальном времени}: через Clash API WebSocket
\item \textbf{Диагностика}: 25 автоматических проверок
\item \textbf{Журнал}: три источника (EclipVPN/sing-box/xray) с фильтрами
\item \textbf{Расписание}: включение и выключение по времени
\item \textbf{Бэкап}: экспорт и импорт настроек
\end{itemize}

\section{Установка}

\subsection{Требования}

\begin{itemize}
\item Роутер с OpenWrt 24.10+ (проверено на Routerich AX3000, aarch64)
\item 256\,МБ ОЗУ, 40\,МБ свободной flash-памяти
\item SSH-доступ как root
\item Подписка на VLESS+Reality серверы (Remnawave или аналог)
\end{itemize}

\subsection{Зависимости}

\begin{lstlisting}
opkg update
opkg install jq coreutils-base64 kmod-nft-tproxy sing-box
\end{lstlisting}

\begin{notebox}
\textbf{Внимание:} \texttt{xray} из opkg — устаревший (25.1.30) и содержит критические баги XHTTP. Установите 26.4.17+ вручную с GitHub-релизов \texttt{XTLS/Xray-core}.
\end{notebox}

\subsection{Установка xray}

\begin{lstlisting}
# На локальной машине:
curl -sL "https://github.com/XTLS/Xray-core/releases/download/v26.4.17/Xray-linux-arm64-v8a.zip" -o /tmp/xray.zip
unzip -o /tmp/xray.zip xray -d /tmp/

# Загрузка на роутер (scp не работает на busybox — используем cat | ssh):
cat /tmp/xray | ssh root@192.168.1.1 "cat > /usr/bin/xray && chmod +x /usr/bin/xray"

# Проверка:
ssh root@192.168.1.1 "/usr/bin/xray version"
\end{lstlisting}

\subsection{Деплой файлов EclipVPN}

\begin{lstlisting}
ssh root@192.168.1.1 "mkdir -p /usr/lib/eclipvpn /www/flowvpn /etc/eclipvpn /tmp/eclipvpn"

for f in usr/bin/eclipvpn usr/lib/eclipvpn/parse.sh usr/lib/eclipvpn/xray.sh \
         usr/lib/eclipvpn/singbox.sh usr/lib/eclipvpn/network.sh \
         www/cgi-bin/eclipvpn www/eclipvpn/index.html etc/init.d/eclipvpn; do
    cat "$f" | ssh root@192.168.1.1 "cat > /$f && chmod +x /$f"
done
\end{lstlisting}

\subsection{Первый запуск}

\begin{lstlisting}
# Указать ссылку на подписку
eclipvpn set-config subscription_url "https://example.com/api/sub/xxx"

# Загрузить серверы и правила маршрутизации
eclipvpn update

# Запустить
eclipvpn start

# Включить автозапуск при загрузке
/etc/init.d/eclipvpn enable
\end{lstlisting}

\section{Веб-интерфейс}

После запуска откройте \texttt{http://192.168.1.1/eclipvpn/} в браузере.

\subsection{Вкладки}

\begin{tabularx}{\textwidth}{l X}
\toprule
\textbf{Вкладка} & \textbf{Описание} \\
\midrule
Маршрут & Главная: переключатель VPN, выбор сервера, режим, скорость \\
Серверы & Список серверов из подписки, пинг, тарифицируемые помечены «Тариф» \\
Устройства & Подключённые клиенты, переключение VPN/Direct для каждого \\
Журнал & Соединения с фильтрами (VPN/Direct/DNS/Ошибки/EclipVPN/xray) \\
Диагностика & 25 автоматических проверок системы \\
Настройки & Подписка, расписание, автообновление, бэкап \\
\bottomrule
\end{tabularx}

\subsection{Режимы маршрутизации}

\begin{tabularx}{\textwidth}{l X}
\toprule
\textbf{Режим} & \textbf{Поведение} \\
\midrule
\texttt{all\_vpn} & Всё через VPN, кроме geosite-ru / geoip-ru. По умолчанию. \\
\texttt{blocklist} & Через VPN только домены из списка Shadowrocket. Остальное — напрямую. \\
\bottomrule
\end{tabularx}

\section{Командная строка}

\subsection{Управление}

\begin{lstlisting}
eclipvpn status                   # JSON со статусом, сервером, IP
eclipvpn start                    # Запуск
eclipvpn stop                     # Остановка с полной очисткой nft+dnsmasq
eclipvpn restart                  # Soft restart — без утечек DNS
eclipvpn update                   # Обновить подписку и правила
\end{lstlisting}

\subsection{Серверы и режим}

\begin{lstlisting}
eclipvpn set-server 0             # Выбрать сервер по индексу
eclipvpn set-mode all_vpn         # Режим маршрутизации
eclipvpn set-mode blocklist
eclipvpn ping                     # Пинг всех серверов
\end{lstlisting}

\subsection{Домены}

\begin{lstlisting}
eclipvpn add-domain youtube.com vpn       # Принудительно через VPN
eclipvpn add-domain rutube.ru direct      # Принудительно напрямую
eclipvpn remove-domain youtube.com vpn
eclipvpn get-domains                      # Список текущих
\end{lstlisting}

\subsection{Диагностика и журнал}

\begin{lstlisting}
eclipvpn diagnose                 # Быстрая диагностика
eclipvpn healthcheck              # 25 проверок
eclipvpn get-log 100 all          # Журнал
eclipvpn clear-log                # Очистить
eclipvpn system-info              # Версии, RAM, flash, uptime
\end{lstlisting}

\subsection{Бэкап}

\begin{lstlisting}
eclipvpn backup > eclipvpn-backup.json
cat eclipvpn-backup.json | eclipvpn restore
\end{lstlisting}

\section{Архитектура}

\begin{lstlisting}
Клиент -> dnsmasq(:53) -> sing-box DNS(:5354) -> dns-proxy(VPN) / dns-direct
       -> nftables tproxy -> sing-box(:2154) -> xray SOCKS(:1080) -> VPN
                                              -> direct (RU)
\end{lstlisting}

\subsection{Ключевые порты}

\begin{tabularx}{\textwidth}{l l X}
\toprule
\textbf{Порт} & \textbf{Процесс} & \textbf{Назначение} \\
\midrule
1080 & xray & SOCKS-прокси для sing-box \\
2154 & sing-box & tproxy inbound \\
5354 & sing-box & DNS inbound (от dnsmasq) \\
9090 & sing-box & Clash API + WebSocket трафика \\
\bottomrule
\end{tabularx}

\subsection{nftables}

\begin{itemize}
\item Таблица: \texttt{table ip eclipvpn} (\textbf{ip}, не \texttt{inet} — \texttt{tproxy} не работает в \texttt{inet})
\item Mark \texttt{0x8000} (32768): обход fw4 (\texttt{meta mark \& 0x00008000 accept})
\item Прерывание: трафик роутера, локальные сети, IP VPN-серверов идут напрямую
\end{itemize}

\subsection{DNS}

\begin{itemize}
\item dnsmasq → \texttt{127.0.0.1\#5354} (sing-box DNS)
\item DNS заблокированных доменов → через VPN (\texttt{tls://1.1.1.1})
\item DNS остальных → напрямую (\texttt{udp://1.1.1.1})
\item DoT (порт 853) и QUIC (UDP 443) блокируются — приложения переключаются на TCP
\end{itemize}

\section{Безопасность}

\subsection{Clash API}

При первом запуске генерируется случайный 32-символьный секрет в \texttt{/etc/eclipvpn/clash\_secret} (\texttt{chmod 600}). UI получает его через \texttt{/status} и использует в WebSocket. Раньше был хардкод \texttt{"eclipvpn"}.

\subsection{CGI-валидация}

Все параметры от UI проходят валидацию на границе CGI:

\begin{tabularx}{\textwidth}{l X}
\toprule
\textbf{Параметр} & \textbf{Допускается} \\
\midrule
\texttt{index} & только цифры \\
\texttt{domain} & \texttt{[a-zA-Z0-9.*\_-]}, без \texttt{..} и без точки в начале \\
\texttt{ip} & валидный IPv4 dotted-quad \\
\texttt{route} & \texttt{vpn} или \texttt{direct} \\
\texttt{list} & \texttt{direct} или \texttt{blocklist} \\
\texttt{mode} & \texttt{[a-zA-Z0-9\_-]} \\
\texttt{on}/\texttt{off} & \texttt{HH:MM} формат \\
\bottomrule
\end{tabularx}

POST body ограничен 64\,КБ. Восстановление бэкапа идёт через stdin (не через argv).

\section{Watchdog}

Каждые 20 секунд проверяет VPN-соединение через текущий SOCKS. После 2 фейлов подряд:

\begin{enumerate}
\item Перебирает следующие серверы из подписки
\item Пропускает тарифицируемых операторов (МТС, Т2, Билайн, Мегафон) — белые списки потеряют смысл при переключении
\item Пробует каждого через probe xray на уникальном порту
\item Первый рабочий — \texttt{soft restart} с новым сервером
\item Спит \texttt{interval}\,×\,3 для стабилизации, чтобы не вызвать каскадный failover
\end{enumerate}

Защита от двух экземпляров через лок-файл \texttt{/tmp/eclipvpn/watchdog.lock}.

\section{Устранение неполадок}

\subsection{VPN не подключается}

\begin{enumerate}
\item \texttt{eclipvpn diagnose} — должно быть «xray работает» и «VPN IP получен»
\item \texttt{eclipvpn ping} — какой сервер отвечает
\item \texttt{logread -e eclipvpn | tail -20} — последние события
\end{enumerate}

\subsection{Сайт не открывается через VPN}

\begin{enumerate}
\item Журнал → проверить идёт ли домен через VPN или Direct
\item Если Direct, но должен быть VPN: добавить домен в «Принудительно через VPN»
\item Если VPN, но не работает: попробовать другой сервер
\end{enumerate}

\subsection{AI пишет «не поддерживается в вашей стране»}

geosite-ru включает \texttt{google.com} → Gemini идёт напрямую с домашнего IP. Решение:
\begin{lstlisting}
eclipvpn add-domain gemini.google.com vpn
eclipvpn add-domain ai.google.com vpn
\end{lstlisting}

\subsection{Приложение пишет «нет сети»}

Кэш DNS и соединений. Авиарежим вкл/выкл на телефоне сбрасывает.

\subsection{Всё сломалось — нет интернета}

\begin{lstlisting}
ssh root@192.168.1.1
eclipvpn stop
# Если не помогает — экстренное восстановление:
nft delete table ip eclipvpn
killall xray sing-box
uci delete dhcp.@dnsmasq[0].server
uci set dhcp.@dnsmasq[0].noresolv='0'
uci commit dhcp
/etc/init.d/dnsmasq restart
\end{lstlisting}

\section{Ресурсы}

\begin{tabularx}{\textwidth}{l X}
\toprule
\textbf{Ресурс} & \textbf{Расход} \\
\midrule
RAM & ${\sim}112$\,МБ (xray 64\,МБ + sing-box 46\,МБ) \\
Flash & ${\sim}33$\,МБ (xray 32{,}5\,МБ + скрипты 112\,КБ) \\
CPU & ${\sim}0\%$ idle, всплески при handshake \\
\bottomrule
\end{tabularx}

\section{Файлы}

\begin{tabularx}{\textwidth}{l X}
\toprule
\textbf{Путь} & \textbf{Назначение} \\
\midrule
\texttt{/usr/bin/eclipvpn} & Главный CLI-скрипт \\
\texttt{/usr/lib/eclipvpn/parse.sh} & Парсер VLESS-подписки \\
\texttt{/usr/lib/eclipvpn/xray.sh} & Генерация xray SOCKS-конфига \\
\texttt{/usr/lib/eclipvpn/singbox.sh} & Генерация sing-box-конфига \\
\texttt{/usr/lib/eclipvpn/network.sh} & nftables tproxy + dnsmasq \\
\texttt{/www/cgi-bin/eclipvpn} & CGI API \\
\texttt{/www/flowvpn/index.html} & Веб-интерфейс \\
\texttt{/etc/eclipvpn/config.json} & Конфигурация \\
\texttt{/etc/eclipvpn/clash\_secret} & Случайный секрет Clash API \\
\texttt{/etc/init.d/eclipvpn} & Автозапуск \\
\texttt{/tmp/eclipvpn/} & Состояние, логи, кэш \\
\bottomrule
\end{tabularx}

\vfill
\begin{center}
\small\color{gray}
GitHub: \url{https://github.com/1eningrad/flowvpn} · MIT License
\end{center}

\end{document}
